\section{Instant-NGP}
Các phương pháp biểu diễn đồ họa dựa trên mạng neural (neural graphics primitives), thường sử dụng Multi-Layer Perceptrons (MLP) để tham số hóa các hàm biểu diễn hình dạng (ví dụ: SDF \cite{Park2019DeepSDF}) hoặc trường bức xạ (ví dụ: NeRF \cite{Mildenhall2020NeRF}), đã đạt được những thành tựu ấn tượng \cite{Muller2022InstantNGP}. Tuy nhiên, việc huấn luyện và tính toán các MLP lớn thường tốn kém tài nguyên và thời gian \cite{Muller2022InstantNGP_source_8}. Các kỹ thuật mã hóa đầu vào (input encoding) như mã hóa tần số (frequency encoding) \cite{Mildenhall2020NeRF, Vaswani2017Attention} giúp cải thiện chất lượng nhưng vẫn đòi hỏi MLP lớn và hội tụ chậm \cite{Muller2022InstantNGP_source_78}. Các phương pháp mã hóa tham số (parametric encoding) sử dụng cấu trúc dữ liệu phụ trợ như lưới dày (dense grid) hoặc cây (tree) \cite{Liu2020NSVF, Takikawa2021NGLOD} có thể tăng tốc huấn luyện bằng cách giảm kích thước MLP, nhưng lại gặp vấn đề về bộ nhớ lớn hoặc yêu cầu các thuật toán cập nhật cấu trúc phức tạp (pruning, merging), không hiệu quả trên GPU hiện đại \cite{Muller2022InstantNGP_source_26, Muller2022InstantNGP_source_28, Muller2022InstantNGP_source_55, Muller2022InstantNGP_source_74, Muller2022InstantNGP_source_81, Muller2022InstantNGP_source_92}.

Để giải quyết những thách thức này, Müller và cộng sự đã đề xuất Instant Neural Graphics Primitives (Instant-NGP), một phương pháp đột phá dựa trên kỹ thuật mã hóa đầu vào mới gọi là \textbf{Multiresolution Hash Encoding} (Mã hóa Băm Đa Phân giải) \cite{Muller2022InstantNGP}. Ý tưởng cốt lõi là thay thế MLP lớn bằng một MLP \textit{nhỏ} kết hợp với một cấu trúc dữ liệu hiệu quả chứa các tham số có thể huấn luyện được \cite{Muller2022InstantNGP_source_9}.

\subsection{Phương pháp Multiresolution Hash Encoding}

Mã hóa băm đa phân giải hoạt động như sau (xem Hình 3 trong \cite{Muller2022InstantNGP}):
\begin{enumerate}
    \item \textbf{Cấu trúc đa mức (Multiresolution Levels):} Sử dụng $L$ (thường là 16) mức phân giải khác nhau. Mỗi mức $l$ tương ứng với một lưới ảo (virtual grid) có độ phân giải $N_l$, tăng dần theo cấp số nhân từ $N_{min}$ đến $N_{max}$ \cite{Muller2022InstantNGP_source_112, Muller2022InstantNGP_source_113, Muller2022InstantNGP_source_114}.
    \item \textbf{Bảng băm đặc trưng (Feature Hash Tables):} Mỗi mức $l$ được liên kết với một bảng (mảng) chứa tối đa $T$ vector đặc trưng (feature vectors) $F$-chiều (thường $F=2$), gọi là $\theta_l$. $T$ là một siêu tham số quan trọng, cân bằng giữa chất lượng, bộ nhớ và hiệu năng \cite{Muller2022InstantNGP_source_105, Muller2022InstantNGP_source_112, Muller2022InstantNGP_source_125}.
    \item \textbf{Tra cứu và Băm (Lookup and Hashing):} Với một tọa độ đầu vào $x \in \mathbb{R}^d$, tại mỗi mức $l$:
        \begin{itemize}
            \item Xác định các đỉnh (corners) $2^d$ của ô lưới ảo (voxel) bao quanh $x$.
            \item Ánh xạ tọa độ nguyên của các đỉnh này tới các chỉ số (indices) trong bảng đặc trưng $\theta_l$:
                \begin{itemize}
                    \item Ở các mức phân giải thô (số đỉnh lưới $\le T$): Ánh xạ 1-1.
                    \item Ở các mức phân giải mịn (số đỉnh lưới $> T$): Sử dụng một hàm băm không gian (spatial hash function), ví dụ $h(\mathbf{x}) = (\bigoplus_{i=1}^d x_i \pi_i) \pmod T$, để lấy chỉ số \cite{Muller2022InstantNGP_source_104, Muller2022InstantNGP_source_117, Muller2022InstantNGP_source_118, Muller2022InstantNGP_source_121, Teschner2003SpatialHashing}.
                \end{itemize}
        \end{itemize}
    \item \textbf{Xử lý Xung đột Băm (Collision Handling):} Instant-NGP \textit{không} sử dụng các cơ chế xử lý xung đột băm tường minh. Thay vào đó, nó dựa vào hai yếu tố:
        \begin{itemize}
            \item Tối ưu hóa dựa trên gradient: Khi nhiều vị trí bị băm vào cùng một mục trong bảng, gradient của chúng sẽ được cộng dồn (hoặc trung bình). Gradient từ các mẫu "quan trọng" hơn (ví dụ: điểm trên bề mặt nhìn thấy trong NeRF) sẽ chiếm ưu thế và điều khiển giá trị của mục đó \cite{Muller2022InstantNGP_source_34, Muller2022InstantNGP_source_150, Muller2022InstantNGP_source_153}.
            \item Mạng MLP theo sau: Mạng neural nhỏ $m$ học cách phân giải các xung đột còn lại bằng cách kết hợp thông tin từ nhiều mức phân giải (các mức thô thường không có xung đột) \cite{Muller2022InstantNGP_source_99, Muller2022InstantNGP_source_119, Muller2022InstantNGP_source_142, Muller2022InstantNGP_source_149}.
        \end{itemize}
    \item \textbf{Nội suy và Kết hợp (Interpolation and Concatenation):} Các vector đặc trưng $F$-chiều được lấy ra từ bảng băm tại các đỉnh voxel được nội suy tuyến tính đa chiều (d-linear interpolation) dựa trên vị trí tương đối của $x$ trong voxel đó \cite{Muller2022InstantNGP_source_105, Muller2022InstantNGP_source_123}. Kết quả nội suy từ tất cả $L$ mức được nối (concatenate) lại với nhau (cùng các đầu vào phụ trợ $\xi$ như hướng nhìn nếu cần) để tạo thành vector đầu vào $y$ cho MLP $m(y; \Phi)$ \cite{Muller2022InstantNGP_source_106, Muller2022InstantNGP_source_124}.
\end{enumerate}

\subsection{Ưu điểm và Cài đặt}

Cách tiếp cận này mang lại nhiều lợi ích đáng kể:
\begin{itemize}
    \item \textbf{Tốc độ huấn luyện cực nhanh:} Cho phép huấn luyện các biểu diễn neural chất lượng cao chỉ trong vài giây đến vài phút trên một GPU duy nhất, nhanh hơn nhiều bậc độ lớn so với các phương pháp trước đó \cite{Muller2022InstantNGP_source_1, Muller2022InstantNGP_source_6, Muller2022InstantNGP_source_16, Muller2022InstantNGP_source_131, Muller2022InstantNGP_source_293, Muller2022InstantNGP_source_345}.
    \item \textbf{Hiệu năng tính toán cao:} Render thời gian thực (vài chục ms ở độ phân giải HD) \cite{Muller2022InstantNGP_source_16, Muller2022InstantNGP_source_255, Muller2022InstantNGP_source_272}.
    \item \textbf{Chất lượng cao:} Đạt được chất lượng tương đương hoặc tốt hơn các phương pháp SOTA như NeRF, mip-NeRF, NSVF, NGLOD trong nhiều tác vụ \cite{Muller2022InstantNGP_source_6, Muller2022InstantNGP_source_30, Muller2022InstantNGP_source_138, Muller2022InstantNGP_source_216, Muller2022InstantNGP_source_292}.
    \item \textbf{Tính linh hoạt và Đơn giản:} Mã hóa không phụ thuộc tác vụ, dễ cấu hình (chủ yếu qua $T$), không cần cơ chế cập nhật cấu trúc dữ liệu phức tạp \cite{Muller2022InstantNGP_source_7, Muller2022InstantNGP_source_29, Muller2022InstantNGP_source_36, Muller2022InstantNGP_source_94, Muller2022InstantNGP_source_195, Muller2022InstantNGP_source_343}. Tự động thích ứng với phân bố dữ liệu \cite{Muller2022InstantNGP_source_158, Muller2022InstantNGP_source_159}.
    \item \textbf{Hiệu quả trên GPU:} Cấu trúc bảng băm dễ song song hóa, tra cứu $O(1)$, tận dụng cache hiệu quả. Toàn bộ hệ thống được cài đặt bằng các CUDA kernel hợp nhất (fully-fused) trong framework \texttt{tiny-cuda-nn} \cite{Muller2021TinyCuDNN, Muller2022InstantNGP_source_15, Muller2022InstantNGP_source_37, Muller2022InstantNGP_source_38, Muller2022InstantNGP_source_39, Muller2022InstantNGP_source_100, Muller2022InstantNGP_source_167}.
\end{itemize}

Instant-NGP sử dụng MLP rất nhỏ (thường chỉ 2 lớp ẩn, 64 neuron) \cite{Muller2022InstantNGP_source_178}, lưu trữ đặc trưng dạng half-precision, tối ưu bằng Adam với $\epsilon$ nhỏ ($10^{-15}$) \cite{Kingma2014Adam, Muller2022InstantNGP_source_168, Muller2022InstantNGP_source_184}. Đối với NeRF, nhóm tác giả sử dụng kiến trúc gồm MLP mật độ và MLP màu, kết hợp occupancy grid để tăng tốc ray marching \cite{Muller2022InstantNGP_source_262, Muller2022InstantNGP_source_270, Muller2022InstantNGP_source_490}.

\subsection{Ứng dụng và Hạn chế}

Instant-NGP đã được chứng minh hiệu quả trên nhiều tác vụ như biểu diễn ảnh gigapixel, học hàm SDF, Neural Radiance Caching \cite{Muller2021NRC}, và đặc biệt là NeRF \cite{Muller2022InstantNGP_source_1}. Nó cho phép huấn luyện NeRF chất lượng cao nhanh hơn 20-60 lần so với frequency encoding trên cùng implementation tối ưu \cite{Muller2022InstantNGP_source_281, Muller2022InstantNGP_source_299}.

Một hạn chế nhỏ là hiện tượng nhiễu hạt (microstructure/graininess) có thể xuất hiện do xung đột băm, đặc biệt rõ trên bề mặt SDF \cite{Muller2022InstantNGP_source_228, Muller2022InstantNGP_source_245, Muller2022InstantNGP_source_328}.

Nhìn chung, Instant-NGP với mã hóa băm đa phân giải là một bước tiến quan trọng, cung cấp một giải pháp thay thế hiệu quả, nhanh chóng và linh hoạt cho các cấu trúc dữ liệu và MLP lớn trong lĩnh vực neural graphics primitives.